\section{\label{sec:defs}Distribution functions}


Throughout this work we focus on flavor non-singlet unpolarized quasi and 
light-front PDFs. The extension to polarized quasi PDFs is straightforward. The 
flavor singlet case introduces additional mixing with the gluon distribution, 
but the principles are similar. We also assume that the
quarks are massless and ignore complications arising from the correct 
treatment of heavy flavors, a subject of continued study for light-front 
PDFs (for reviews, see, for example, 
\cite{Thorne:2008xf,Olness:2008px,Binoth:2010nha}).

\subsection{Bare PDFs}

In the following section, when we use the term ``bare'' we mean finite
matrix elements determined with some regulator at finite cutoff. We leave the 
regulator implicit in this discussion, although one can have in mind 
dimensional regularization if desired. These bare matrix elements require 
renormalization in some scheme before one can remove the regulator (or, on the 
lattice, take the continuum limit). This usage follows that of the extensive 
discussions of light-front PDFs in, for example, 
\cite{Collins:1981uw,Collins:2011zzd}.

We denote bare light-front PDFs by 
$f^{(0)}(\xi)$. Light-front PDFs are frequently represented by 
$f_{j/N}^{(0)}(\xi)$, where $j$ 
denotes the quark flavor and $N$ the nucleon 
species, but here we will be considering only non-singlet distributions, for 
which we can neglect mixing between parton species, and work with sufficient 
generality that the nucleon species is not relevant to our discussion. We 
use light-front coordinates, 
$(x^+,x^-,\boldsymbol{x}_\mathrm{T})$ such that $x^\pm = (t\pm z)/\sqrt{2}$, 
and define $\xi=k^+/P^+$. We use $\xi$ to distinguish this variable from the 
Bjorken-$x$ parameter that characterizes the kinematics of scattering 
experiments and is given in 
terms of the experimental momentum transfer $Q^2=-q^2$ and hadron momentum $P$ 
by $x=Q^2/(2P\cdot q)$. The bare PDF is defined as \cite{Collins:2011zzd}
\begin{equation}
f^{(0)}(\xi) = \int_{-\infty}^\infty 
\frac{\mathrm{d}\omega^-}{4\pi}
e^{-i\xi P^+\omega^-} \left\langle P \left| 
T\,\overline{\psi}(0,\omega^-,\mathbf{0}_\mathrm{T})
W(\omega^-,0)\gamma^+\frac{\lambda^a}{2}\psi(0) \right| 
P\right\rangle_\mathrm{C}.
\end{equation}
Here $T$ is the 
time-ordering operator, $\psi$ is a quark 
field, and the subscript C indicates that the vacuum 
expectation value has been subtracted (in other words, only connected 
contributions are included). The operator $W(\omega^-,0)$ is the Wilson line,
\begin{equation}
W(\omega^-,0) = 
{\cal P}\exp\left[-ig_0\int_0^{\omega^-}\mathrm{d}y^-A^+_\alpha(0,y^-,
\mathbf{0}_{\mathrm{T}})T_\alpha\right],
\end{equation}
with ${\cal P}$ the path-ordering operator, $g_0$ the QCD bare coupling, 
and $A^\mu = A^\mu_\alpha T_\alpha$ the $SU(3)$ gauge potential with generator 
$T_\alpha$ (summation over color index $\alpha$ is implicit). The target 
state, $|P\rangle$, is a spin-averaged, exact momentum eigenstate with 
relativistic normalization
\begin{equation}
\langle P' | P \rangle = 
(2\pi)^32P^+\delta\left(P^+-P^{\prime\,+}\right)\delta^{(2)}
\left(\mathbf{P}_\mathrm{T} - \mathbf{P}'_\mathrm{T}\right).
\end{equation}

We define the moments of bare PDFs as
\begin{equation}
a_{0}^{(n)} = \int_0^1 \mathrm{d}\xi\,\xi^{n-1}\left[
f^{(0)}(\xi)+(-1)^n\overline{f}^{(0)}(\xi)\right] = \int_{-1}^1 
\mathrm{d}\xi\,\xi^{n-1}
f(\xi),
\end{equation}
where $\overline{f}^{(0)}(\xi)$ is the anti-quark PDF and the second equality 
follows from the relation of the quark to 
anti-quark PDFs
\begin{equation}\label{eq:pdf2apdf}
f^{(0)}(-\xi) = -\overline{f}^{(0)}(\xi),
\end{equation}
which holds for the bare distributions if the quark and anti-quarks fields are 
classical, or quantized using light-front quantization \cite{Collins:1981uw}.

We can relate these bare moments, $a_{0}^{(n)}$, to matrix elements of 
twist-two operators via
\begin{equation}\label{eq:opePDF}
\left\langle P | {\cal O}_{0}^{\{\mu_1\ldots \mu_n\}} | P 
\right\rangle = 2a_{0}^{(n)} \left(P^{\mu_1}\cdots P^{\mu_n} - 
\mathrm{traces}\right).
\end{equation}
Here the bare
twist-two operators are 
\begin{equation}\label{eq:baretwist2def}
{\cal 
O}_0^{\{\mu_1\cdots \mu_n\}} = i^{n-1}\overline{\psi}(0)
\gamma^{\{\mu_1}D^{\mu_2}\cdots 
D^{\mu_n\}}\frac{\lambda^a}{2}\psi(0)-\operatorname{traces}.
\end{equation}
In these expressions the braces denote symmetrization, $D^\mu$ is the 
symmetric covariant 
derivative, $\lambda^a$ are $SU(2)$ flavor matrices, and the subtraction of the 
trace terms ensures that the operator 
transforms irreducibly under $SU(2)_{\mathrm{L}}\otimes SU(2)_{\mathrm{R}}$. 

\subsection{Renormalized PDFs}

To this point we have considered the bare light-front PDFs, with the 
understanding that such objects are evaluated with some regulator that renders 
the bare distributions finite. We now introduce
renormalized  light-front PDFs. We stress that in this section we 
consider a renormalization scheme that respects rotational symmetry and, for 
definiteness, one can have in mind the $\ms$ scheme.  Complications will arise 
if  a regulator that breaks rotational invariance, such as 
the lattice regulator, is used. We do not discuss such complications here, 
because we will avoid explicit computations of moments 
at finite lattice spacing. All correlation functions computed on the lattice 
can be renormalized and extrapolated to the continuum limit, provided that  
no power divergent mixing exists. In the next section, we propose a smeared 
correlation function that does not have power-divergent mixing.

In general, renormalized light-front PDFs are written in terms of a 
kernel, ${\cal Z}(\zeta/\xi,\mu)$, as
\begin{equation}
f(\xi,\mu) = \int_\xi^1 \frac{\mathrm{d}\zeta}{\zeta}{\cal 
Z}\left(\frac{\zeta}{\xi},\mu\right)f^{(0)}(\zeta),
\end{equation}
where $\mu$ is some renormalization 
scale. We do not need to consider mixing between parton species for 
non-singlet distributions. In terms of the renormalized light-front PDF, the 
renormalized Mellin 
moments are
\begin{equation}
a^{(n)}(\mu) = \int_0^1 \mathrm{d}\xi\,\xi^{n-1}\left[
f(\xi,\mu)+(-1)^n\overline{f}(\xi,\mu)\right] = \int_{-1}^1 
\mathrm{d}\xi\,\xi^{n-1}
f(\xi,\mu),
\end{equation}
which can be related to matrix elements of renormalized twist-two operators, 
${\cal O}^{\{\nu_1\ldots \nu_n\}}(\mu) =Z_{\cal O}(\mu){\cal 
O}_0^{\{\nu_1\ldots \nu_n\}}$, via
\begin{equation}\label{eq:RopePDF}
\left\langle P | {\cal O}^{\{\nu_1\ldots \nu_n\}}(\mu) | P 
\right\rangle = 2a^{(n)}(\mu) \left(P^{\nu_1}\cdots P^{\nu_n} - 
\mathrm{traces}\right).
\end{equation}
This relation holds provided the light-front PDFs and twist-two operators are 
renormalized in the same scheme \cite{Collins:1981uw}.



\subsection{Smeared  quasi PDFs}

We construct a finite quasi PDF matrix element by smearing both the fermion and 
gauge fields via the gradient flow 
\cite{Narayanan:2006rf,Luscher:2011bx,Luscher:2013cpa}. The gradient flow is a 
deterministic 
evolution of the original quark and gluon fields in a new dimension, 
generally referred to as the ``flow time'', towards a classical minimum of the 
QCD action \cite{Luscher:2011bx,Luscher:2013cpa}. The flow-time evolution is 
chosen to remove ultraviolet fluctuations, which corresponds to smearing out 
the quark and gluon fields in real space, with a smearing scale that is 
proportional to the square-root of the flow time. Here we will not describe in 
detail the gradient flow, but refer the reader to the recent reviews 
\cite{Luscher:2013vga,Ramos:2015dla} for more details and applications. 

For our purposes, it is sufficient that the gradient flow has the following 
properties. First, the gradient flow serves as a gauge-invariant ultraviolet 
regulator. Second, given a renormalized theory at zero flow time, the matrix 
elements of smeared fields are automatically finite, up to a multiplicative 
wave-function renormalization for the fermion fields \cite{Luscher:2013cpa}, 
which can be removed by introducing ringed fermion fields 
\cite{Makino:2014taa,Hieda:2016lly}. 
Third, the lattice matrix elements of smeared fields
remain finite in the continuum limit, provided the flow time is fixed in 
physical units \cite{Luscher:2013cpa,Monahan:2015lha}. In essence, the gradient 
flow allows 
one to replace the lattice regulator with a new smearing-scale regulator. This 
last fact allows one to determine the continuum limit of lattice matrix 
elements 
of, for example, twist-two operators, without power-divergent mixing. In the 
continuum, because the gradient flow respects rotational symmetry, the mixing 
between twist-two operators is then reduced to ordinary mixing with 
coefficients 
that depend on the smearing scale and not powers of the inverse lattice spacing 
\cite{Monahan:2015lha}.

We denote the ringed fermion fields at flow time $\tau$ by
$\overline{\chi}(x;\tau)$ and $\chi(x;\tau)$, and the corresponding Wilson line 
at the same flow time, constructed from the smeared gauge fields 
$B_\mu(x;\tau)$, by ${\bf\cal W}(0,z;\tau)$.
We start with the matrix element
\begin{equation}\label{eq:hdef}
h^{(s)}\left(\frac{z}{\sqrt{\tau}},\sqrt{\tau}P_z,
\sqrt{\tau}\Lambda_{\mathrm{QCD}},\sqrt{\tau}M_\mathrm{N}\right) = 
\frac{1}{2P_z}\left\langle 
P_z \left| \overline{\chi}(z;\tau) 
{\bf\cal W}(0,z;\tau)\gamma_z\frac{\lambda^a}{2}\chi(0;\tau)\right| 
P_z\right\rangle_\mathrm{C},
\end{equation}
which, being dimensionless, depends only on dimensionless combinations of 
scales. We note that the flow time has units of length-squared. The subscript C 
indicates that disconnected contributions to this 
matrix element have been removed. The ringed fermion fields require no wave
function renormalization and this smeared matrix element is finite 
provided the flow time, $\tau$, is non-zero and fixed in physical units, 
because correlation functions constructed from smeared fields are finite 
\cite{Luscher:2011bx,Luscher:2013cpa}. Note that divergences will appear in the 
limit of vanishing flow time and the matrix element will then require
renormalization.

We then define the quasi PDF 
\cite{Ji:2013dva,Ji:2014gla} as
\begin{equation}\label{eq:qpdfdef}
q^{\,(s)}\left(\xi,\sqrt{\tau}P_z,\sqrt{\tau}\Lambda_{\mathrm{QCD}}, 
\sqrt{\tau}M_\mathrm{N}\right) 
= 
\int_{-\infty}^\infty \frac{\mathrm{d}z}{2\pi} e^{i\xi z 
P_z} P_z\, 
h^{(s)}(\sqrt{\tau}z,\sqrt{\tau}P_z,\sqrt{\tau}\Lambda_{\mathrm{QCD}}, 
\sqrt{\tau}M_\mathrm{N}),
\end{equation}
where $\xi$ is a dimensionless parameter that can be naively interpreted as the 
longitudinal momentum fraction of the parton in the nucleon $N$.  This 
interpretation is not correct in Euclidean space, however, and instead $\xi$ 
should be viewed as a dimensionless momentum variable in a Fourier 
transformation.  

In practice, the smeared matrix element $h$ is determined from lattice 
computations at finite lattice spacing, $a$, as 
\begin{equation}\label{eq:hacont}
h^{(s)}\left(\frac{z}{\sqrt{\tau}},\sqrt{\tau}P_z,\sqrt{\tau}\Lambda_{\mathrm{
QCD}},
\sqrt { \tau }
M_\mathrm{N}\right) =  \lim_{a\to 0}
h\left(\frac{z}{a},\frac{\sqrt{\tau}}{a},aP_z,a\Lambda_{\mathrm{QCD}},aM_\mathrm
{N}
\right) ,
\end{equation}
where $\sqrt{\tau}$ and $P_z$ are held fixed and
\begin{align}
h\Big(\frac{z}{a},\frac{\sqrt{\tau}}{a},aP_z,a\Lambda_{\mathrm{QCD}},{} & 
aM_\mathrm{N}\Big) = \nonumber\\
{} & \frac{1}{2aP_z}
\left\langle 
aP_z \left| 
\overline{\chi}\left(\frac{z}{a};\frac{\sqrt{\tau}}{a}\right) 
W\left(0,\frac{z}{a};\frac{\sqrt{\tau}}{a}\right)\gamma_z\frac{\lambda^a}{2}
\chi\left(0;\frac { \sqrt{\tau}}{a} \right)\right| 
aP_z\right\rangle_\mathrm{C}.
\end{align}

